\section{Basic Puzzle Properties}
\label{sec:PuzzleLemmas}

This section establishes elementary properties of the puzzle pieces defined in Section \ref{sec:Puzzle}.

\subsection{Nested Structure}

The puzzle pieces form a nested sequence:
\[ D_0(z) \supseteq D_1(z) \supseteq D_2(z) \supseteq \dots \]
This follows immediately from the definition involving sublevel sets of the Green's function: if $G_c(w) < 2^{-(n+1)}$, then certainly $G_c(w) < 2^{-n}$.

\textbf{Theorem} (\texttt{dynamical\_puzzle\_piece\_nested}):
The dynamical puzzle piece of depth $n+1$ is contained in the piece of depth $n$.

\begin{lstlisting}[language=Lean]
theorem dynamical_puzzle_piece_nested (c : $\mathbb{C}$) (n : $\mathbb{N}$) :
    DynamicalPuzzlePiece c (n + 1) 0 $\subseteq$ DynamicalPuzzlePiece c n 0
\end{lstlisting}

\subsection{Self-Containment}

By definition, the point $z$ is always contained in its puzzle piece $D_n(z)$ (assuming the Green's function value allows it, or if $z=0$ and $c \in \mathcal{M}$, then $G_c(0)=0 < 2^{-n}$).

\textbf{Lemma} (\texttt{mem\_dynamical\_puzzle\_piece\_self}):
If $c \in \mathcal{M}$, then $0$ is contained in $D_n(0)$ for all $n$.

\begin{lstlisting}[language=Lean]
lemma mem_dynamical_puzzle_piece_self (c : $\mathbb{C}$) (hc : c $\in$ MandelbrotSet) (n : $\mathbb{N}$) :
    0 $\in$ DynamicalPuzzlePiece c n 0
\end{lstlisting}

\subsection{Intersection Properties}

We also relate the puzzle pieces to the filled Julia set $K_c$. Since $K_c = \{z \mid G_c(z) = 0\}$, and puzzle pieces are defined by $G_c(z) < \epsilon$, we have $K_c \subseteq D_n(0)$ (for the component containing 0).
